\snsection{INTRODUCCIÓN}
\justifying
La electrónica de potencia combina la energía, la electrónica y el control. Este último se encarga del régimen permanente y de las características dinámicas de los sistemas de lazo cerrado. Por otro lado, la energía tiene que ver con el equipo de potencia estática y rotativa o giratoria, para la generación, transmisión y distribución de la misma. Y por último, la electrónica se ocupa de los dispositivos y circuitos de estado sólido requeridos en el procesamiento de señales para cumplir con los objetivos de control deseados. Por ende, la electrónica de potencia puede ser definida como ``la aplicación de la electrónica de estado sólido para el control y la conversión de la energía eléctrica''.

%Esta rama de la electrónica se basa, en primer término, en la conmutación de dispositivos semiconductores de potencia como, transistores, tiristores, diodos, entre otros. 
Con el desarrollo de la tecnología de los semiconductores de potencia, las capacidades del manejo de la energía y la velocidad de conmutación de los dispositivos de potencia ha mejorado rápidamente. El desarrollo de la tecnología de los microprocesadores-microcomputadoras tiene un gran impacto sobre el control y la síntesis de la estrategia de control para los dispositivos semiconductores de potencia. Además, al trabajar con frecuencias mayores, permite la reducción del tamaño de los componentes, \parencite{RASHID}.

%El equipo de electrónica de potencia moderno utiliza (1) semiconductores de potencia, que pueden compararse con el músculo, y (2) microelectrónica, que tiene el poder y la inteligencia del cerebro. 

Por otra parte, el equipamiento didáctico permite a los estudiantes experimentar y entender mejor los principios teóricos mediante la práctica. Estos equipos facilitan el aprendizaje activo, permitiendo a los estudiantes realizar experimentos, observar fenómenos y analizar resultados en un entorno controlado. Esta experiencia práctica es vital para el desarrollo de habilidades críticas, resolver problemas complejos y prepararse para los desafíos del mundo real. 

El proyecto consiste en el diseño y construcción de un sistema didáctico modular con el objetivo de facilitar la enseñanza y el aprendizaje de la electrónica de potencia. Este sistema permitirá a estudiantes o entusiastas experimentar con  diversos módulos, tales como rectificadores, inversores, convertidores DC-DC, entre otros. Además, mediante la observación de señales de entrada y salida y el ajuste de parámetros en tiempo real, se permite experimentar en un entorno seguro los distintos módulos y como se comportarían en una situación de aplicación real.

Además, contará con un manual de usuario que explique el funcionamiento teórico de cada módulo y una guía práctica que oriente sobre los parámetros a modificar y los efectos a analizar. Con este enfoque, se espera que los usuarios desarrollen una comprensión profunda y tangible de los principios de la materia.

El diseño modular del sistema permitirá que, en el futuro, se puedan incorporar expansiones o módulos personalizados según necesidades del usuario o avances tecnológicos. De este modo, el proyecto queda abierto a mejoras continuas y a la adaptación a nuevas áreas de estudio de la electrónica.

El sistema propuesto surge de la necesidad de contar con herramientas para el aprendizaje de la electrónica de potencia que sean accesibles, seguras y efectivas. Desarrollar un sistema didáctico modular permitirá la enseñanza y el aprendizaje de los conceptos de la cátedra, sin necesidad de contar con un laboratorio especializado y costoso. Asimismo, se podrán visualizar de manera clara y simple los efectos de cambiar parámetros en los distintos circuitos, agilizando los tiempos de aprendizaje, en especial debido al poco tiempo disponible en el año que se tiene para la enseñanza de la materia, teniendo en cuenta el amplio contenido que se debe abordar. 

\begin{comment}
\justifying
Escriba aquí una introducción al tema.
Algunos autores denominan Introducción como “Antecedentes” (no confundir con el
marco teórico) o como “Resumen”. Por lo tanto, la introducción debe concentrar, con fluidez y precisión, de manera discursiva, los principales elementos del proyecto, permitiendo al lector familiarizarse con él.

La redacción de la introducción debe ser ligera y amena, una especie de diálogo que debe motivar al lector a continuar leyendo el anteproyecto. Al terminar de leer la introducción, el lector:

Entenderá el proceso de motivación y decisión de llevar a cabo el proyecto; se ubicará en el contexto y el enfoque desde el cual el investigador abordará el tema;

Contará con información preliminar para comprender y evaluar el proyecto, sin tener que consultar otros documentos para clarificarlo.

No debe haber un corte abrupto entre la introducción y la descripción, es decir esta última debe ser una continuidad de la introducción.

La introducción tiene la tarea o finalidad de presentar el primer acercamiento al lector, ha de ser suficientemente amplia como para dar una idea general, pero precisa, para señalar el área temática y el problema de investigación, sin dejar de lado otros puntos a contemplar.

Diferentes fuentes de referencia en lo que respecta a la metodología pueden proporcionar muy buenas ideas de los contenidos esperados, uno de los posibles listados de estos puede ser:

\begin{itemize}
    \item Presentación del área de conocimiento, indicando el contexto en que se inserta el proyecto
    \item Área temática del trabajo, pues la extensión de un área es en extremo amplia
    \item Problema de investigación, ya el señalamiento concreto de cuál es el problema que se identifica y una breve presentación de lo que lo caracteriza como problema. Puede encontrarse aquí también la enunciación del objetivo general del trabajo, redactado de manera llana y preliminar
    \item Señalamientos de las teorías y corrientes de pensamiento empleadas para el abordaje del problema y el cumplimiento del objetivo
    \item Indicaciones de la metodología del trabajo que se presenta
    \item Especificación de la estructura de presentación de los contenidos del escrito 
\end{itemize}

Tal como puede pensarse, este listado es breve e incluso simple, pero el orden de los elementos puede presentar y permitir variaciones, igualmente, la extensión esperada de una introducción no es mucha, solo la suficiente para presentar los puntos y despertar el interés del lector.

Esta última idea es relevante, se escribe pensando en el lector, por tanto se ha de cuidar el enlace de las ideas, la argumentación, la presentación de los elementos que le sean necesarios para abordar el tema y entender el enfoque del proyecto. En tal sentido, es bueno recordar una comparación que se usa en estos casos:

La introducción es al Informe lo que un trailer cinematográfico es a su película.

\end{comment}